\documentclass[11pt,a4paper]{article}
\usepackage{fontspec}
\setmainfont{TeX Gyre Pagella}
\newfontfamily\headfont{Inter}[UprightFont=*-ExtraBold,BoldFont=*-ExtraBold,Extension=.otf,Path=/usr/share/fonts/opentype/inter/]
\usepackage[a4paper,margin=2.3cm]{geometry}
\usepackage{graphicx,booktabs,float,caption,array,xcolor}
\usepackage[hidelinks]{hyperref}
\usepackage{titlesec}
\usepackage[round]{natbib}
\usepackage{setspace}
\setstretch{1.05}
\renewcommand{\topfraction}{0.9}
\renewcommand{\bottomfraction}{0.8}
\renewcommand{\textfraction}{0.1}
\renewcommand{\floatpagefraction}{0.75}
\setlength{\parindent}{0pt}
\setlength{\parskip}{0.55em}
\titleformat{\section}{\headfont\Large}{\thesection}{0.8em}{}
\titleformat{\subsection}{\headfont\large}{\thesubsection}{0.8em}{}
\captionsetup{font=small,labelfont=bf}
\urlstyle{same}
\hypersetup{pdftitle={Life on the Moving Map: 540 Million Years of Animal Fossils on the Drifting Continents},pdfauthor={Tristan Hotham}}

\begin{document}
\begin{center}
{\Large\bfseries Life on the Moving Map\par}
\vspace{0.4em}
{\large 540 million years of animal fossils on the drifting continents,\\from the Paleobiology Database and the PALEOMAP plate model\par}
\vspace{1.2em}
{\large Tristan Hotham\par}
\vspace{0.3em}
{\small Independent researcher\\
ORCID: 0000-0003-4931-9349\\
%%DOILINE%%\par}
\vspace{0.8em}
2026
\end{center}

\begin{quote}
\small
\begin{center}\textbf{Abstract}\end{center}
Every fossil has two locations: where it is found today and where the animal lived, which plate tectonics has since moved. I place {{records.n|int}} fossil occurrences of {{records.n_groups|int}} animal groups, from trilobites to hominids, from the Paleobiology Database at their ancient positions, computed with the PALEOMAP plate model, and draw them on PALEOMAP continental outlines for {{frames.n|int}} ages from 540 million years ago to today. I describe the assembly and break-up of Pangaea with a simple index, the share of all land in the largest connected land mass, which exceeds 85\% from {{pangaea.first_over_85|int}} to {{pangaea.last_over_85|int}} million years ago and peaks at {{pangaea.peak_pct|f1}}\% at {{pangaea.peak_age|int}} million years ago (today {{pangaea.today_pct|f1}}\%). Counting distinct genera in 5 million year bins, trilobites make up {{leaders.Paleozoic.pct|pc}} of Paleozoic genera in these groups, insects {{leaders.Mesozoic.pct|pc}} of Mesozoic and mammals {{leaders.Cenozoic.pct|pc}} of Cenozoic genera. Dinosaur genera from the Late Triassic and Early Jurassic, when Pangaea was whole, are found on more than one present-day continent more than twice as often ({{dino.multi_pangaea.pct|f1}}\%) as Cretaceous genera ({{dino.multi_cretaceous.pct|f1}}\%). A simple last-appearance measure picks out the end-Permian ({{extinctions.2.lost_pct|f1}}\% of genera last seen) and end-Cretaceous extinctions but not the end-Triassic. Throughout, the record is strongly shaped by where palaeontologists work: {{sampling.na_eu_pct|pc}} of occurrences come from North America and Europe, which hold {{sampling.na_eu_land_pct|pc}} of the land. An interactive map lets readers move through time, switch any fossil between its ancient and present-day position, and search for genera.

\textbf{Keywords:} palaeogeography; fossil record; Paleobiology Database; PALEOMAP; Pangaea; dinosaurs; mass extinctions; sampling bias; data visualisation
\end{quote}

\section{Introduction}
Continents move a few centimetres a year. Over the 540 million years since animals with hard parts became common, that is enough to carry a piece of crust from the southern tropics to the Arctic, to gather nearly all land into one supercontinent, Pangaea, and to break it apart again. Fossils travel with the rocks that hold them. A trilobite found in Wales, a coal-swamp insect from northern England and an ichthyosaur from the Dorset coast all lie at about 52 degrees north today, but the median collection of each age lay at {{drift.uk_420|f1}}, {{drift.uk_300|f1}} and {{drift.uk_190|f1}} degrees of palaeo-latitude respectively (collections from the United Kingdom dated 420 to 440, 300 to 320 and 190 to 210 million years ago).

The Paleobiology Database \citep{peters2016} records both positions for most collections: the present-day coordinates, and palaeo-coordinates computed with a plate model. This paper uses those palaeo-coordinates, from the PALEOMAP model of \citet{scotese2016}, together with PALEOMAP continental outlines reconstructed through the GPlates Web Service \citep{muller2018}, to build an animated map of the Earth through time on which each fossil sits where it lived. It is a descriptive and visual study, meant to make familiar facts (the drifting continents, Pangaea, the succession of dominant groups, the mass extinctions) visible in one place from open data, and to be honest about how much of that picture is shaped by the record rather than by life.

\section{Data}
\textbf{Fossil occurrences.} I downloaded all occurrences of {{records.n_queries|int}} named taxa from the Paleobiology Database data service (version 1.2), data accessed 2026, requesting palaeo-coordinates from the PALEOMAP model at each collection's midpoint age ({{records.rows_raw|int}} rows; Table~\ref{tab:groups}). An occurrence is the record of a taxon in a collection, a set of fossils from one place and rock unit. The taxa were chosen to match the animals people tend to ask about: trilobites, sea scorpions, armoured fish (placoderms), ammonoids, sharks and rays, insects, temnospondyl amphibians, non-mammal synapsids, turtles, ichthyosaurs, plesiosaurs, crocodylomorphs, pterosaurs, the three main dinosaur groups, birds, mosasaurs, mammals, and three mammal groups (whales, proboscideans and hominids). The database is licensed CC BY 4.0.

\textbf{Continental outlines.} PALEOMAP coastline and continental block polygons \citep{scotese2016,scotese2018} were reconstructed with the GPlates Web Service (model PALEOMAP) every 10 million years from 0 to 540 million years ago and at eleven further ages, including the five mass extinctions. They come from the same plate model as the fossil palaeo-coordinates, so the two line up. These polygons are present-day coastlines and blocks carried back in time, not palaeo-shorelines.

\textbf{Present-day places.} Natural Earth 1:10m countries and states and provinces (public domain) give the present-day country and province of each collection.

{{TABLE_GROUPS}}

\section{Methods}
\subsection{Cleaning and grouping}
The query results overlap: Dinosauria includes the three dinosaur subgroups and the birds, and Mammalia includes whales, proboscideans and hominids. Each occurrence is kept once, in its most specific group, removing {{records.dupes|int}} repeats. Three query results reach beyond the animals intended and were corrected: the name Proboscidea also matches an insect taxon ({{records.pro_not_mammal|int}} records outside class Mammalia removed; the insects are in the Insecta query); the database classifies the living amphibians (Lissamphibia) inside Temnospondyli, so {{records.tem_to_lis|int}} records of frogs, salamanders, caecilians and albanerpetontids form a separate group; and {{records.the_to_ave|int}} records of bird orders returned under non-avian theropods were moved to the birds. {{records.dino_late|int}} non-bird dinosaur records with midpoint ages after 66 million years ago were dropped. This gives {{records.n_groups|int}} groups.

\subsection{Ages and positions}
Each occurrence takes the midpoint of its collection's age range. Occurrences whose age range exceeds {{records.max_range|int}} million years ({{records.wide_age|int}}) are excluded from all analyses. Of {{records.paleo_missing|int}} occurrences without a PALEOMAP position, {{records.paleo_filled|int}} younger than 1 million years were given their present position; the remaining {{records.dropped_no_paleo|int}} were dropped. The analysis set has {{records.n|int}} occurrences in {{records.collections|int}} collections and {{records.genera|int}} genera, from {{records.countries|int}} present-day countries; {{records.with_genus_pct|f1}}\% of occurrences are identified to genus.

\subsection{Map frames}
For each of the {{frames.n|int}} outline ages, the map shows occurrences with midpoints within 5 million years of the age, except at the five extinction ages (445, 372, 252, 201 and 66 million years ago), where it shows the 10 million years before the extinction. Occurrences are binned by group into 1 degree cells of palaeo-latitude and longitude; each cell keeps its number of occurrences and collections, mean palaeo-position and mean present-day position, present-day place, age span and taxon names. The outlines are merged into land masses, simplified for the web and drawn in the Equal Earth projection.

\subsection{Pangaea index}
For each age I merged the outline polygons, closed gaps narrower than about 0.3 degrees (blocks of one continent are separated by hairline gaps in the source), joined pieces across the 180th meridian, and measured each land mass on the sphere in a cylindrical equal-area projection. The index is the area of the largest land mass divided by the total land area. It depends on the gap threshold and on the outlines' fidelity, and measures connection, not shared climate or fauna.

\subsection{Diversity, latitude and extinction}
Diversity is the number of distinct genera with midpoints in each 5 million year bin, with no correction for sampling. Because midpoints of stage-length ranges cluster, bins fill unevenly; one bin ({{empty_bins.0|int}} to 230 million years ago) is empty. For latitude I take one point per collection and group, and summarise palaeo-latitude every 10 million years by the median and the 10th and 90th percentiles, where there are at least 5 collections. For the extinctions, a genus with at least two collections is counted as standing at a boundary if its oldest midpoint is older than the boundary and its youngest is younger than 10 million years before it; it is lost if its youngest midpoint falls in those 10 million years. The same measure at every other 10 million year step gives a background (median {{ext_bg_median|f1}}\%, 90th percentile {{ext_bg_p90|f1}}\%, {{ext_bg_n|int}} steps). For dinosaurs, a genus is counted as found on more than one continent if its collections in an interval lie on two or more present-day continents.

\section{Results}
\subsection{The moving map}
Figure~\ref{fig:multiples} shows six frames. At 500 million years ago the record is almost entirely trilobites, in shallow seas around the scattered continents. By 300 million years ago insects, ammonoids and temnospondyls dominate the frame, and the continents are converging. At 200 million years ago nearly all land is one block and the finds cluster in western North America, Europe and southern Africa. By 66 million years ago the Atlantic and Indian Oceans are open and the record is full of ammonites, sharks, turtles and dinosaurs. The present frame is dominated by mammals and birds. Figure~\ref{fig:two} shows the 300 million year frame at ancient and present positions side by side.

\begin{figure}[tbp]\centering\includegraphics[width=\textwidth]{f_multiples.png}
\caption{Fossil occurrences at six ages, at their palaeo-positions on the PALEOMAP outlines, coloured by group. Each dot is a 1 degree cell; size grows with the number of occurrences.}\label{fig:multiples}\end{figure}

\begin{figure}[tbp]\centering\includegraphics[width=\textwidth]{f_twoaddresses.png}
\caption{The same occurrences of 295 to 305 million years ago at their palaeo-positions (left) and present-day positions (right).}\label{fig:two}\end{figure}

\subsection{Pangaea}
The Pangaea index (Figure~\ref{fig:pangaea}) stays between 56 and 71\% from 540 to 280 million years ago, when Gondwana was the largest land mass. It rises in steps through the Carboniferous and Permian as the continents collide, reaching {{pangaea.at.270|f1}}\% at 270 million years ago, stays above 85\% for {{pangaea.years_over_85|int}} million years and peaks at {{pangaea.peak_pct|f1}}\% at {{pangaea.peak_age|int}} million years ago. It then falls quickly: {{pangaea.at.170|f1}}\% at 170, {{pangaea.at.150|f1}}\% at 150 and {{pangaea.at.100|f1}}\% at 100 million years ago. Today the largest land mass, Africa with Eurasia, holds {{pangaea.today_pct|f1}}\% of the land.

\begin{figure}[tbp]\centering\includegraphics[width=.92\textwidth]{f_pangaea.png}
\caption{Pangaea index: share of all land in the largest connected land mass. Shaded: above 85\%. Dashed: the five mass extinctions.}\label{fig:pangaea}\end{figure}

\subsection{Which groups dominated the record}
Figure~\ref{fig:lanes} scales each group's genus counts to its own peak. The succession is familiar: trilobites and sea scorpions in the early Paleozoic, armoured fish in the Devonian, temnospondyls and synapsids in the late Paleozoic and Triassic, marine reptiles and dinosaurs in the Mesozoic, and mammals, whales, proboscideans and hominids in the Cenozoic. As a share of all genera in these groups, trilobites lead the Paleozoic ({{leaders.Paleozoic.pct|f1}}\%), insects the Mesozoic ({{leaders.Mesozoic.pct|f1}}\%) and mammals the Cenozoic ({{leaders.Cenozoic.pct|f1}}\%). Insects have the most genera of any group ({{groups.ins.genera|int}}), and their record is highly concentrated: the 50 richest of {{facts.ins_coll|int}} insect collections hold {{facts.ins_top50_coll_pct|f1}}\% of insect occurrences.

\begin{figure}[tbp]\centering\includegraphics[width=.95\textwidth]{f_lanes.png}
\caption{Distinct genera per 5 million years for each group, each scaled to its own maximum. Groups in order of first appearance; dashed lines: mass extinctions.}\label{fig:lanes}\end{figure}

\subsection{Dinosaurs}
The data hold {{dino.n|int}} occurrences of non-bird dinosaurs in {{dino.genera|int}} genera: {{dino.by_sub.the.genera|int}} theropod, {{dino.by_sub.sau.genera|int}} sauropodomorph and {{dino.by_sub.orn.genera|int}} ornithischian genera. Genus counts are low through the Triassic and Jurassic and highest in the Late Cretaceous (Figure~\ref{fig:dinos}); the richest bin, {{dino.peak_bin.1|int}} to {{dino.peak_bin.0|int}} million years ago, has {{dino.peak_genera|int}} genera, {{dino.peak_na_pct|f1}}\% of its occurrences from North America. Of {{dino.multi_pangaea.genera|int}} genera from the Late Triassic and Early Jurassic, {{dino.multi_pangaea.pct|f1}}\% occur on two or more present-day continents, against {{dino.multi_cretaceous.pct|f1}}\% of {{dino.multi_cretaceous.genera|int}} Cretaceous genera (Figure~\ref{fig:cosmo}). The Late Jurassic stands out at {{dino.intervals.3.multi_pct|f1}}\%, while the North Atlantic was still narrow. All tyrannosaurid and ceratopsid occurrences are from North America, Asia or Europe; {{facts.abe_south_india_pct|f1}}\% of abelisaurid occurrences are from the southern continents and India. {{dino.polar_coll|int}} dinosaur collections lie beyond 60 degrees of palaeo-latitude and {{dino.polar75_coll|int}} beyond 75 degrees, chiefly in Alaska and Victoria, Australia.

\begin{figure}[tbp]\centering\includegraphics[width=.92\textwidth]{f_dinos.png}
\caption{Dinosaur genera per 5 million years, stacked by group, with birds as a dashed line.}\label{fig:dinos}\end{figure}

\begin{figure}[tbp]\centering\includegraphics[width=\textwidth]{f_cosmo.png}
\caption{Top: non-bird dinosaur occurrences at 210, 150 and 70 million years ago. Bottom: share of genera in each interval found on two or more present-day continents (bars) and the Pangaea index at the nearest outline age (line).}\label{fig:cosmo}\end{figure}

\subsection{Latitude}
Median absolute palaeo-latitudes range from {{lat.hom|f1}} degrees for hominids to {{lat.syn|f1}} degrees for non-mammal synapsids (Figure~\ref{fig:lat}). Trilobite collections have a median of {{lat.tri|f1}} degrees, ammonoids {{lat.amm|f1}} and theropods {{lat.the|f1}}. Because most collections come from the Northern Hemisphere continents, which have been studied most, latitudinal spreads describe the record rather than the full ranges of the animals.

\begin{figure}[tbp]\centering\includegraphics[width=\textwidth]{f_latitude.png}
\caption{Palaeo-latitude of collections by 10 million year step for six groups: median (line) and 10th to 90th percentiles (band). Shaded: the tropics.}\label{fig:lat}\end{figure}

\subsection{Mass extinctions}
Figure~\ref{fig:ext} shows the last-appearance measure at the five extinctions. The end-Permian is the highest ({{extinctions.2.lost_pct|f1}}\% of {{extinctions.2.standing|int}} genera), above the 90th percentile of the background. The end-Ordovician ({{extinctions.0.lost_pct|f1}}\%), Late Devonian ({{extinctions.1.lost_pct|f1}}\%) and end-Cretaceous ({{extinctions.4.lost_pct|f1}}\%) are well above the background median; the end-Triassic ({{extinctions.3.lost_pct|f1}}\%) is not, probably because much of its toll fell on groups not included here and because Late Triassic sampling is thin. At the end-Cretaceous every genus of non-bird dinosaur, plesiosaur and mosasaur with at least two collections is last seen before the boundary, as are {{extinctions.4.by_group.amm.lost_pct|f1}}\% of ammonoids, against {{extinctions.4.by_group.ins.lost_pct|f1}}\% of insects.

\begin{figure}[tbp]\centering\includegraphics[width=\textwidth]{f_extinctions.png}
\caption{Share of standing genera (with at least two collections) last seen in the 10 million years before each extinction, by group (groups with at least 10 standing genera; number standing in brackets). Dashed: median of the same measure at other 10 million year steps.}\label{fig:ext}\end{figure}

\subsection{Sampling}
North America holds {{sampling.by_continent.0.pct|f1}}\% of occurrences and Europe {{sampling.by_continent.1.pct|f1}}\%; together {{sampling.na_eu_pct|f1}}\%, against {{sampling.na_eu_land_pct|f1}}\% of the land area. The United States alone has {{sampling.usa_pct|f1}}\%; Africa {{sampling.africa_pct|f1}}\% and South America {{sampling.sam_pct|f1}}\% (Figure~\ref{fig:samp}). The North American and European share is similar in each era (Paleozoic {{sampling.na_eu_by_era.Paleozoic|f1}}\%, Mesozoic {{sampling.na_eu_by_era.Mesozoic|f1}}\%, Cenozoic {{sampling.na_eu_by_era.Cenozoic|f1}}\%) but swings widely between 5 million year bins.

\begin{figure}[tbp]\centering\includegraphics[width=\textwidth]{f_sampling.png}
\caption{Left: occurrences by present-day continent. Right: share of occurrences from North America and Europe per 5 million years (bins with at least 20 collections); dashed, their share of today's land.}\label{fig:samp}\end{figure}

\section{Discussion}
Placed at their ancient positions, open fossil data reproduce the broad textbook picture of the Phanerozoic: the dominance of trilobites in early Paleozoic seas, the assembly and break-up of Pangaea, the provincial faunas of the Cretaceous, and the turnover at the great extinctions. That such a picture emerges from simple counts is a tribute to the decades of work collected in the Paleobiology Database. But each pattern is also a pattern of sampling. Genus counts follow the area of exposed rock of each age and the effort spent on it \citep{raup1972,smith2007,alroy2008}; the high Late Cretaceous dinosaur counts reflect, in part, the badlands of western North America; and the latitudinal patterns reflect the positions of the well-studied Northern Hemisphere continents. The cosmopolitanism result fits the expectation that faunas became more provincial as Pangaea broke up \citep{brusatte2015}, but genus definitions and national traditions of naming add noise that this simple measure cannot remove.

\section{Limitations}
First, occurrence counts are not abundance, and genus counts are not true diversity: no correction for sampling is applied. Second, ages are midpoints of often long ranges, so the timing of events is blurred by several million years, and bins fill unevenly. Third, palaeo-positions come from one plate model at each collection's own age; other models differ, uncertainty grows with age, and map dots can sit up to 5 million years away from the outline age. Fourth, the outlines are present-day coastlines and blocks moved in time, not ancient shorelines; flooded continental interiors are drawn as land. Fifth, the chosen groups are a small and uneven part of animal life, and the query corrections, although checked, may leave some misassigned records. Sixth, I could not consult PALEOMAP's own reuse terms from the build environment; the outlines are credited in full and shared with that caveat.

\section{Conclusion}
A map of the moving Earth with its fossils at their ancient positions is a vivid way to see deep time, and it can be built entirely from open data. It shows Pangaea forming and breaking apart, the changing cast of animals, and the great extinctions, and it shows, just as clearly, that the fossil record is a map of where rock is exposed and where people have looked.

\section*{Data and code availability}
The interactive map is at \url{{{meta.page}}}. Data and code are at \url{{{meta.repo}}}%%ARCHIVE%%. Fossil data: The Paleobiology Database (paleobiodb.org), data accessed 2026, CC BY 4.0. Derived data are released under CC BY 4.0 and code under the MIT licence.

\section*{AI statement}
I wrote the text, with drafting and editing assistance from Claude (Anthropic), which also helped prepare the analysis code. I reviewed and checked all text and results.

\begin{thebibliography}{99}
\bibitem[Alroy et al.(2008)]{alroy2008} Alroy, J. et al. (2008) Phanerozoic trends in the global diversity of marine invertebrates. \emph{Science}, 321, 97 to 100.
\bibitem[Brusatte et al.(2015)]{brusatte2015} Brusatte, S.L. et al. (2015) The extinction of the dinosaurs. \emph{Biological Reviews}, 90, 628 to 642.
\bibitem[Cohen et al.(2013)]{ics} Cohen, K.M., Finney, S.C., Gibbard, P.L. and Fan, J.-X. (2013, updated) The ICS International Chronostratigraphic Chart. \emph{Episodes}, 36, 199 to 204.
\bibitem[Müller et al.(2018)]{muller2018} Müller, R.D. et al. (2018) GPlates: building a virtual Earth through deep time. \emph{Geochemistry, Geophysics, Geosystems}, 19(7), 2243 to 2259. GPlates Web Service: \url{https://gws.gplates.org/}
\bibitem[Natural Earth(2026)]{ne} Natural Earth (2026) Admin 0 countries and admin 1 states and provinces, 1:10m. Public domain. \url{https://www.naturalearthdata.com/}
\bibitem[Peters and McClennen(2016)]{peters2016} Peters, S.E. and McClennen, M. (2016) The Paleobiology Database application programming interface. \emph{Paleobiology}, 42, 1 to 7. Data: The Paleobiology Database (paleobiodb.org), data accessed 2026, CC BY 4.0.
\bibitem[Raup(1972)]{raup1972} Raup, D.M. (1972) Taxonomic diversity during the Phanerozoic. \emph{Science}, 177, 1065 to 1071.
\bibitem[Scotese(2016)]{scotese2016} Scotese, C.R. (2016) PALEOMAP PaleoAtlas for GPlates and the PaleoData Plotter Program. PALEOMAP Project.
\bibitem[Scotese and Wright(2018)]{scotese2018} Scotese, C.R. and Wright, N. (2018) PALEOMAP Paleodigital Elevation Models (PaleoDEMS) for the Phanerozoic. PALEOMAP Project.
\bibitem[Smith and McGowan(2007)]{smith2007} Smith, A.B. and McGowan, A.J. (2007) The shape of the Phanerozoic marine palaeodiversity curve: how much can be predicted from the sedimentary rock record of western Europe? \emph{Palaeontology}, 50, 765 to 774.
\end{thebibliography}
\end{document}
